\bookchapter{Git Every Day}{ch:08-git-every-day}

\begin{chapterintro}

This chapter turns the ideas from Chapter 4 into commands. You'll start a repository, learn where your changes live before and after a commit, read history and differences, keep unwanted files out, work on branches, share with a remote, and undo the most common mistakes. These are the Git commands you'll use every working day.

\end{chapterintro}

\emph{The problem.} I want to save my work properly, not as \texttt{final\_}\allowbreak{}\texttt{v2\_}\allowbreak{}\texttt{really\_}\allowbreak{}\texttt{final.}\allowbreak{}\texttt{py}.

\emph{The question.} What are the daily Git commands, and what does each one do to my project?

The examples build a tiny project, \texttt{recipes}, that scales ingredient amounts. You need the name and email you set in the Chapter 4 lab.

\codeneed{11}

\section{Starting a Repository}\label{08-git-every-day:starting-a-repository}

Any folder becomes a Git repository with \texttt{git\ init}:

\codeneed{5}\begin{codeblock}

\begin{Shaded}
\begin{Highlighting}[]
\FunctionTok{mkdir} \AttributeTok{{-}p}\NormalTok{ \textasciitilde{}/projects/recipes}
\BuiltInTok{cd}\NormalTok{ \textasciitilde{}/projects/recipes}
\FunctionTok{git}\NormalTok{ init}
\end{Highlighting}
\end{Shaded}

\end{codeblock}

\codeneed{1}\begin{codeblock}
\begin{Verbatim}[formatcom=\outputstyle]
Initialized empty Git repository in /home/ada/projects/recipes/.git/
\end{Verbatim}
\end{codeblock}

Git created one hidden folder, \texttt{.git}. That's the repository: every commit and branch lives inside it. Your files stay where they are. Delete \texttt{.git} and the folder is an ordinary folder again, with all its history gone.

\codeneed{10}

\texttt{git\ status} is the command you'll run most. It tells you where you are and what has changed:

\codeneed{7}\begin{codeblock}

\begin{Shaded}
\begin{Highlighting}[]
\FunctionTok{git}\NormalTok{ status}
\end{Highlighting}
\end{Shaded}

\end{codeblock}

\codeneed{5}\begin{codeblock}
\begin{Verbatim}[formatcom=\outputstyle]
On branch main

No commits yet

nothing to commit (create/copy files and use "git add" to track)
\end{Verbatim}
\end{codeblock}

\section{Three Places Your Work Lives}\label{08-git-every-day:three-places-your-work-lives}

Git keeps your work in three places, and every command moves changes between them. Think of posting a parcel. Your desk is where you make things. The box is where you put exactly what you want to send. The post office stamps the box and keeps a record forever.

\begin{itemize}
\tightlist
\item
  The \textbf{working tree}\index{working tree} is your desk: the normal files you edit.
\item
  The \textbf{staging area}\index{staging area} is the box: the changes you've chosen for the next commit. (Git also calls it the index.)
\item
  The repository is the post office: the commits, stored safely in \texttt{.git}.
\end{itemize}

\codeneed{6}\begin{codeblock}
\begin{Verbatim}[breaklines=false, fontsize=\fontsize{8.8}{8.68}\selectfont, fontfamily=diagrammono, formatcom=\diagramlines]
  working tree      staging area       repository
  (your files)      (next commit)      (history)
       │                 │                 │
       │ ── git add ───► │                 │
       │                 │ ─ git commit ─► │
       │ ◄──────── git restore ─────────── │
\end{Verbatim}
\end{codeblock}

Why the box? Because you often change several things at once, and they belong in separate commits. Staging lets you choose.

\codeneed{10}

\section{The Everyday Loop: Add and Commit}\label{08-git-every-day:the-everyday-loop-add-and-commit}

Create \texttt{scale.py} in your editor:

\codeneed{4}\begin{codeblock}

\begin{Shaded}
\begin{Highlighting}[]
\KeywordTok{def}\NormalTok{ scale(amount, factor):}
    \ControlFlowTok{return}\NormalTok{ amount }\OperatorTok{*}\NormalTok{ factor}

\BuiltInTok{print}\NormalTok{(scale(}\DecValTok{200}\NormalTok{, }\DecValTok{2}\NormalTok{))}
\end{Highlighting}
\end{Shaded}

\end{codeblock}

\codeneed{9}

Git notices the new file but doesn't track it yet: \texttt{git\ status} now lists \texttt{scale.py} under ``Untracked files''. \texttt{git\ add} puts it in the box; \texttt{git\ commit} records the snapshot. The \texttt{-m} gives the message:

\codeneed{6}\begin{codeblock}

\begin{Shaded}
\begin{Highlighting}[]
\FunctionTok{git}\NormalTok{ add scale.py}
\FunctionTok{git}\NormalTok{ commit }\AttributeTok{{-}m} \StringTok{"Add function to scale a recipe"}
\end{Highlighting}
\end{Shaded}

\end{codeblock}

\codeneed{3}\begin{codeblock}
\begin{Verbatim}[formatcom=\outputstyle]
[main (root-commit) 4978b71] Add function to scale a recipe
 1 file changed, 4 insertions(+)
 create mode 100644 scale.py
\end{Verbatim}
\end{codeblock}

\texttt{4978b71} is the start of the new commit's ID. \texttt{git\ log} shows the history: each commit's full ID, author, date, and message.

That's the loop you'll repeat all day: edit, \texttt{git\ status}, \texttt{git\ add}, \texttt{git\ commit}.

\begin{callout}{tip}

\texttt{git\ add\ .} stages every change in the current folder. It's handy, but run \texttt{git\ status} first so you don't commit something you didn't mean to.

\end{callout}

\codeneed{11}

\section{Seeing What Changed}\label{08-git-every-day:seeing-what-changed}

Now change \texttt{scale.py} so it rounds results, and add a second example:

\codeneed{5}\begin{codeblock}

\begin{Shaded}
\begin{Highlighting}[]
\KeywordTok{def}\NormalTok{ scale(amount, factor):}
    \ControlFlowTok{return} \BuiltInTok{round}\NormalTok{(amount }\OperatorTok{*}\NormalTok{ factor, }\DecValTok{1}\NormalTok{)}

\BuiltInTok{print}\NormalTok{(scale(}\DecValTok{200}\NormalTok{, }\DecValTok{2}\NormalTok{))}
\BuiltInTok{print}\NormalTok{(scale(}\DecValTok{150}\NormalTok{, }\FloatTok{0.333}\NormalTok{))}
\end{Highlighting}
\end{Shaded}

\end{codeblock}

\codeneed{14}

\texttt{git\ diff} shows changes in the working tree that aren't staged yet. Lines starting with \texttt{-} were removed; lines starting with \texttt{+} were added:

\codeneed{11}\begin{codeblock}

\begin{Shaded}
\begin{Highlighting}[]
\FunctionTok{git}\NormalTok{ diff}
\end{Highlighting}
\end{Shaded}

\end{codeblock}

\codeneed{9}\begin{codeblock}
\begin{Verbatim}[formatcom=\outputstyle]
diff --git a/scale.py b/scale.py
…
@@ -1,4 +1,5 @@
 def scale(amount, factor):
-    return amount * factor
+    return round(amount * factor, 1)
 
 print(scale(200, 2))
+print(scale(150, 0.333))
\end{Verbatim}
\end{codeblock}

\codeneed{11}

Once you \texttt{git\ add} the file, plain \texttt{git\ diff} shows nothing, because nothing is left unstaged. \texttt{git\ diff\ -\/-staged} shows what's in the box instead. Always look before you commit:

\codeneed{8}\begin{codeblock}

\begin{Shaded}
\begin{Highlighting}[]
\FunctionTok{git}\NormalTok{ add scale.py}
\FunctionTok{git}\NormalTok{ commit }\AttributeTok{{-}m} \StringTok{"Round scaled amounts to one decimal place"}
\FunctionTok{git}\NormalTok{ log }\AttributeTok{{-}{-}oneline}
\end{Highlighting}
\end{Shaded}

\end{codeblock}

\codeneed{4}\begin{codeblock}
\begin{Verbatim}[formatcom=\outputstyle]
[main 450c3a8] Round scaled amounts to one decimal place
 1 file changed, 2 insertions(+), 1 deletion(-)
450c3a8 Round scaled amounts to one decimal place
4978b71 Add function to scale a recipe
\end{Verbatim}
\end{codeblock}

\texttt{git\ log\ -\/-oneline} gives one line per commit, newest first. \texttt{git\ show\ 450c3a8} shows a single commit and its changes.

\section{Keeping Files Out with .gitignore}\label{08-git-every-day:keeping-files-out-with-gitignore}

Some files must never be committed: secrets such as \texttt{.env}, and files your tools generate, such as Python's \texttt{\_\_pycache\_\_} folder or a \texttt{.venv}. Git lists them as untracked, which is noise at best and a leak at worst:

\codeneed{4}\begin{codeblock}

\begin{Shaded}
\begin{Highlighting}[]
\FunctionTok{git}\NormalTok{ status }\AttributeTok{{-}{-}short}
\end{Highlighting}
\end{Shaded}

\end{codeblock}

\codeneed{2}\begin{codeblock}
\begin{Verbatim}[formatcom=\outputstyle]
?? .env
?? __pycache__/
\end{Verbatim}
\end{codeblock}

\codeneed{6}

List the patterns to ignore in a file called \textbf{.gitignore}\index{.gitignore} at the top of the project, one per line:

\codeneed{3}\begin{codeblock}
\begin{Verbatim}[formatcom=\outputstyle]
.env
__pycache__/
.venv/
\end{Verbatim}
\end{codeblock}

Now \texttt{git\ status\ -\/-short} shows only \texttt{??\ .gitignore}. Commit the \texttt{.gitignore} itself, so everyone who clones the project ignores the same things.

\begin{callout}{warning}

\texttt{.gitignore} only stops files from being added. If a secret was already committed, ignoring it now doesn't remove it from history. Chapter 12 explains what to do.

\end{callout}

\section{Writing Good Commit Messages}\label{08-git-every-day:writing-good-commit-messages}

A commit message is a note to the person reading the history later, usually you. Write the first line as a short command that completes the sentence ``If applied, this commit will\ldots{}''. Keep it under about 50 characters. If more explanation helps, leave a blank line and write why you made the change.

\Needspace{7\baselineskip}

\begin{booktable}
\begin{longtable}{>{\raggedright\arraybackslash}p{\dimexpr 0.2742\linewidth-1.4516\tabcolsep\relax}>{\raggedright\arraybackslash}p{\dimexpr 0.7258\linewidth-0.5484\tabcolsep\relax}}
\tabletop
\rowcolor{tablehead}\thead{Weak} & \thead{Better} \\
\tablemid
\endfirsthead
\tabletop
\rowcolor{tablehead}\thead{Weak} & \thead{Better} \\
\tablemid
\endhead
\tablebottom
\endlastfoot
\texttt{fix} & \texttt{Fix\ crash\ when\ amount\ is\ zero} \\
\texttt{changes} & \texttt{Round\ scaled\ amounts\ to\ one\ decimal\ place} \\
\texttt{stuff\ for\ Sam} & \texttt{Add\ ounces-}\allowbreak{}\texttt{to-}\allowbreak{}\texttt{grams\ conversion} \\
\end{longtable}
\end{booktable}

One commit should hold one logical change. If your message needs ``and'', consider two commits.

\codeneed{12}

\section{Branches and Merging}\label{08-git-every-day:branches-and-merging}

To try something without touching \texttt{main}, create a branch and switch to it in one step with \texttt{git\ switch\ -c}. \texttt{git\ branch} lists branches; the \texttt{*} marks where HEAD is:

\codeneed{6}\begin{codeblock}

\begin{Shaded}
\begin{Highlighting}[]
\FunctionTok{git}\NormalTok{ switch }\AttributeTok{{-}c}\NormalTok{ metric}
\FunctionTok{git}\NormalTok{ branch}
\end{Highlighting}
\end{Shaded}

\end{codeblock}

\codeneed{3}\begin{codeblock}
\begin{Verbatim}[formatcom=\outputstyle]
Switched to a new branch 'metric'
  main
* metric
\end{Verbatim}
\end{codeblock}

Add a \texttt{to\_grams} function to \texttt{scale.py} and commit it on the branch. (\texttt{git\ commit\ -am} stages every changed tracked file and commits in one step.) Then switch back: \texttt{main} doesn't have the new function yet, because the commit lives only on \texttt{metric}:

\codeneed{11}\begin{codeblock}

\begin{Shaded}
\begin{Highlighting}[]
\FunctionTok{git}\NormalTok{ commit }\AttributeTok{{-}am} \StringTok{"Add ounces{-}to{-}grams conversion"}
\FunctionTok{git}\NormalTok{ switch main}
\FunctionTok{git}\NormalTok{ merge metric}
\end{Highlighting}
\end{Shaded}

\end{codeblock}

\codeneed{7}\begin{codeblock}
\begin{Verbatim}[formatcom=\outputstyle]
[metric 015db0b] Add ounces-to-grams conversion
 1 file changed, 4 insertions(+)
Switched to branch 'main'
Updating f88f44f..015db0b
Fast-forward
 scale.py | 4 ++++
 1 file changed, 4 insertions(+)
\end{Verbatim}
\end{codeblock}

This was a \textbf{fast-forward}\index{fast-forward} merge. \texttt{main} hadn't moved since the branch split, so Git simply slid the \texttt{main} bookmark forward to the branch's commit. No new commit was needed. Delete the finished branch with \texttt{git\ branch\ -d\ metric}.

\codeneed{9}

When both branches have new commits, Git makes a merge commit with two parents. Git opens your editor with a ready-made message; save and close it to finish. \texttt{git\ log\ -\/-graph} draws the result:

\codeneed{6}\begin{codeblock}
\begin{Verbatim}[formatcom=\outputstyle]
*   0ea3aaa Merge branch 'readme'
|\  
| * 97e1215 Add a README
* | 8ea3115 Print a conversion example
|/  
* 015db0b Add ounces-to-grams conversion
\end{Verbatim}
\end{codeblock}

\begin{callout}{tip}

If Git opens an editor you don't know how to leave, you're probably in Vim: type \texttt{:wq} and press Enter. To use a friendlier editor, run \texttt{git\ config\ -}\allowbreak{}\texttt{-}\allowbreak{}\texttt{global\ core.}\allowbreak{}\texttt{editor\ "nano"} (or \texttt{"code\ -\/-wait"} for VS Code).

\end{callout}

\section{Push and Pull}\label{08-git-every-day:push-and-pull}

A remote can be any Git repository you can reach, including a folder on your own disk. A \textbf{bare repository}\index{bare repository} has history but no working tree: nobody edits files in it, it only receives and hands out commits. That's what hosting services like GitHub keep on their servers, and you can make one yourself:

\codeneed{8}\begin{codeblock}

\begin{Shaded}
\begin{Highlighting}[]
\FunctionTok{git}\NormalTok{ init }\AttributeTok{{-}{-}bare}\NormalTok{ \textasciitilde{}/remotes/recipes.git}
\FunctionTok{git}\NormalTok{ remote add origin \textasciitilde{}/remotes/recipes.git}
\FunctionTok{git}\NormalTok{ push }\AttributeTok{{-}u}\NormalTok{ origin main}
\end{Highlighting}
\end{Shaded}

\end{codeblock}

\codeneed{4}\begin{codeblock}
\begin{Verbatim}[formatcom=\outputstyle]
Initialized empty Git repository in /home/ada/remotes/recipes.git/
To /home/ada/remotes/recipes.git
 * [new branch]      main -> main
branch 'main' set up to track 'origin/main'.
\end{Verbatim}
\end{codeblock}

\texttt{git\ remote\ add} gives the remote a short name; origin is the usual name for the main one. \texttt{git\ push\ -u\ origin\ main} sends your \texttt{main} branch there, and \texttt{-u} remembers the pairing, so later you can type just \texttt{git\ push} and \texttt{git\ pull}.

\codeneed{11}

Anyone (including you, on another machine) can now \texttt{git\ clone\ \textasciitilde{}/}\allowbreak{}\texttt{remotes/}\allowbreak{}\texttt{recipes.}\allowbreak{}\texttt{git}, commit, and \texttt{git\ push}. To bring their commits into your copy, pull:

\codeneed{8}\begin{codeblock}

\begin{Shaded}
\begin{Highlighting}[]
\FunctionTok{git}\NormalTok{ pull}
\end{Highlighting}
\end{Shaded}

\end{codeblock}

\codeneed{6}\begin{codeblock}
\begin{Verbatim}[formatcom=\outputstyle]
From /home/ada/remotes/recipes
   0ea3aaa..1d39f76  main       -> origin/main
Updating 0ea3aaa..1d39f76
Fast-forward
 README.md | 1 +
 1 file changed, 1 insertion(+)
\end{Verbatim}
\end{codeblock}

Chapter 11 covers what happens when two people push at once.

\Needspace{17\baselineskip}

\section{Undoing Mistakes}\label{08-git-every-day:undoing-mistakes}

Git can undo almost anything, as long as you pick the right tool for where the mistake lives:

\begin{booktable}
\begin{longtable}{>{\raggedright\arraybackslash}p{\dimexpr 0.6304\linewidth-0.7391\tabcolsep\relax}>{\raggedright\arraybackslash}p{\dimexpr 0.3696\linewidth-1.2609\tabcolsep\relax}}
\tabletop
\rowcolor{tablehead}\thead{Situation} & \thead{Command} \\
\tablemid
\endfirsthead
\tabletop
\rowcolor{tablehead}\thead{Situation} & \thead{Command} \\
\tablemid
\endhead
\tablebottom
\endlastfoot
Throw away edits to a file (not yet staged) & \texttt{git\ restore\ scale.py} \\
Take a file out of the box, keeping the edits & \texttt{git\ restore\ -}\allowbreak{}\texttt{-}\allowbreak{}\texttt{staged\ notes.}\allowbreak{}\texttt{txt} \\
Fix the last commit's message, or add a forgotten file & \texttt{git\ commit\ -\/-amend} \\
Undo a commit that's already pushed & \texttt{git\ revert\ \textless{}id\textgreater{}} \\
Undo the last commit, keeping its changes staged & \texttt{git\ reset\ -\/-soft\ HEAD\textasciitilde{}1} \\
Throw away the last commit and its changes & \texttt{git\ reset\ -\/-hard\ HEAD\textasciitilde{}1} \\
\end{longtable}
\end{booktable}

\codeneed{5}

\texttt{HEAD\textasciitilde{}1} means ``one commit before where I am''. A typo'd message is fixed like this:

\codeneed{2}\begin{codeblock}

\begin{Shaded}
\begin{Highlighting}[]
\FunctionTok{git}\NormalTok{ commit }\AttributeTok{{-}am} \StringTok{"Add serving sise"}
\FunctionTok{git}\NormalTok{ commit }\AttributeTok{{-}{-}amend} \AttributeTok{{-}m} \StringTok{"Add serving size"}
\end{Highlighting}
\end{Shaded}

\end{codeblock}

\codeneed{11}

\texttt{git\ revert} is the safe undo for shared history. Instead of deleting a commit, it adds a new commit that does the opposite, so nobody who already pulled is confused:

\codeneed{8}\begin{codeblock}

\begin{Shaded}
\begin{Highlighting}[]
\FunctionTok{git}\NormalTok{ revert }\AttributeTok{{-}{-}no{-}edit}\NormalTok{ HEAD}
\FunctionTok{git}\NormalTok{ log }\AttributeTok{{-}{-}oneline} \AttributeTok{{-}3}
\end{Highlighting}
\end{Shaded}

\end{codeblock}

\codeneed{5}\begin{codeblock}
\begin{Verbatim}[formatcom=\outputstyle]
[main 2f795a7] Revert "Add serving size"
…
2f795a7 Revert "Add serving size"
ebebf8c Add serving size
1d39f76 Describe the project
\end{Verbatim}
\end{codeblock}

\texttt{git\ reset} moves the branch bookmark backwards, as if later commits never happened. Use it only on commits you haven't pushed.

\begin{callout}{warning}

\texttt{git\ restore\ \textless{}file\textgreater{}} and \texttt{git\ reset\ -\/-hard} discard uncommitted work for good. Git never saw those edits, so it can't give them back. Commit often.

\end{callout}

And if you reset too far? Suppose you ran \texttt{git\ reset\ -\/-hard\ HEAD\textasciitilde{}1} twice, throwing away both the revert and ``Add serving size''. Git keeps a diary of everywhere HEAD has been, the \textbf{reflog}\index{reflog}, and even ``lost'' commits appear there for weeks:

\codeneed{6}\begin{codeblock}

\begin{Shaded}
\begin{Highlighting}[]
\FunctionTok{git}\NormalTok{ reflog }\AttributeTok{{-}4}
\end{Highlighting}
\end{Shaded}

\end{codeblock}

\codeneed{4}\begin{codeblock}
\begin{Verbatim}[formatcom=\outputstyle]
1d39f76 HEAD@{0}: reset: moving to HEAD~1
ebebf8c HEAD@{1}: reset: moving to HEAD~1
2f795a7 HEAD@{2}: revert: Revert "Add serving size"
ebebf8c HEAD@{3}: commit (amend): Add serving size
\end{Verbatim}
\end{codeblock}

\texttt{git\ reset\ -}\allowbreak{}\texttt{-}\allowbreak{}\texttt{hard\ ebebf8c} would bring that commit back. Chapter 10 practises exactly this rescue.

\section{Lab: Save Your Work Properly}\label{08-git-every-day:lab-save-your-work-properly}

Make a new folder and go into it: \texttt{mkdir\ -}\allowbreak{}\texttt{p\ \textasciitilde{}/}\allowbreak{}\texttt{projects/}\allowbreak{}\texttt{lab8}, then \texttt{cd\ \textasciitilde{}/}\allowbreak{}\texttt{projects/}\allowbreak{}\texttt{lab8}.

\begin{enumerate}
\def\labelenumi{\arabic{enumi}.}
\tightlist
\item
  \texttt{git\ init}, then \texttt{git\ status}. \textbf{Expected:} \texttt{On\ branch\ main} and \texttt{No\ commits\ yet}.
\item
  Create \texttt{hello.py} containing \texttt{print("hello")}. Run \texttt{git\ status\ -\/-short}. \textbf{Expected:} \texttt{??\ hello.py}.
\item
  \texttt{git\ add\ hello.py}, then \texttt{git\ status\ -\/-short}. \textbf{Expected:} \texttt{A\ \ hello.py}.
\item
  \texttt{git\ commit\ -}\allowbreak{}\texttt{m\ "Add\ greeting\ script"}. \textbf{Expected:} a line starting \texttt{{[}main\ (root-commit)} and \texttt{1\ file\ changed}.
\item
  Change the text to \texttt{"hello,\ world"} and run \texttt{git\ diff}. \textbf{Expected:} one \texttt{-} line and one \texttt{+} line.
\item
  \texttt{git\ commit\ -}\allowbreak{}\texttt{am\ "Greet\ the\ world"}, then \texttt{git\ log\ -\/-oneline}. \textbf{Expected:} two commits, newest first.
\item
  Create \texttt{.env} and a \texttt{.gitignore} containing \texttt{.env}. Run \texttt{git\ status\ -\/-short}. \textbf{Expected:} only \texttt{??\ .gitignore}. Commit it with \texttt{git\ add\ .gitignore} and \texttt{git\ commit\ -}\allowbreak{}\texttt{m\ "Ignore\ .}\allowbreak{}\texttt{env"}.
\item
  \texttt{git\ switch\ -c\ shout}, change the text to \texttt{"HELLO,\ WORLD"}, commit, \texttt{git\ switch\ main}, \texttt{git\ merge\ shout}. \textbf{Expected:} \texttt{Fast-forward}, and \texttt{python3\ hello.py} prints \texttt{HELLO,\ WORLD}.
\item
  Make a bare remote with \texttt{git\ init\ -}\allowbreak{}\texttt{-}\allowbreak{}\texttt{bare\ \textasciitilde{}/}\allowbreak{}\texttt{remotes/}\allowbreak{}\texttt{lab8.}\allowbreak{}\texttt{git}, then \texttt{git\ remote\ add\ origin\ \textasciitilde{}/}\allowbreak{}\texttt{remotes/}\allowbreak{}\texttt{lab8.}\allowbreak{}\texttt{git} and \texttt{git\ push\ -u\ origin\ main}. \textbf{Expected:} \texttt{*\ {[}new\ branch{]}\ \ \ \ \ \ main\ -}\allowbreak{}\texttt{\textgreater{}\ main}.
\item
  Edit \texttt{hello.py} badly, then \texttt{git\ restore\ hello.py}. \textbf{Expected:} \texttt{git\ status} says \texttt{nothing\ to\ commit,}\allowbreak{}\texttt{\ working\ tree\ clean}.
\end{enumerate}

\section{Summary, Key Terms, and Review Questions}\label{08-git-every-day:summary-key-terms-and-review-questions}

\subsection*{Summary}\label{08-git-every-day:summary}

\begin{itemize}
\tightlist
\item
  \texttt{git\ init} creates the repository (\texttt{.git}). \texttt{git\ status} is your compass.
\item
  Work moves from the working tree, to the staging area (\texttt{git\ add}), to the repository (\texttt{git\ commit}).
\item
  \texttt{git\ diff}, \texttt{git\ diff\ -\/-staged}, \texttt{git\ log}, and \texttt{git\ show} let you see changes and history.
\item
  \texttt{.gitignore} keeps secrets and generated files out. Good messages say what the commit does, in a short command.
\item
  \texttt{git\ switch\ -c} creates a branch; \texttt{git\ merge} joins it back, by fast-forward or with a merge commit.
\item
  \texttt{git\ push} and \texttt{git\ pull} sync with a remote such as \texttt{origin}.
\item
  Undo with \texttt{restore}, \texttt{-\/-amend}, \texttt{revert} (safe for shared history), or \texttt{reset} (local only). The reflog remembers where you've been.
\end{itemize}

\Needspace{13\baselineskip}

\subsection*{Key Terms}\label{08-git-every-day:key-terms}

\begin{booktable}
\begin{longtable}{>{\raggedright\arraybackslash}p{\dimexpr 0.2289\linewidth-1.5422\tabcolsep\relax}>{\raggedright\arraybackslash}p{\dimexpr 0.7711\linewidth-0.4578\tabcolsep\relax}}
\tabletop
\rowcolor{tablehead}\thead{Term} & \thead{Meaning} \\
\tablemid
\endfirsthead
\tabletop
\rowcolor{tablehead}\thead{Term} & \thead{Meaning} \\
\tablemid
\endhead
\tablebottom
\endlastfoot
Working tree & The project files you see and edit \\
Staging area & The changes chosen for the next commit (also called the index) \\
.gitignore & A file listing patterns Git should not track \\
Fast-forward & A merge that just moves the branch pointer forward \\
Bare repository & A repository with history but no working tree, used as a remote \\
Reflog & Git's record of everywhere HEAD has pointed \\
\end{longtable}
\end{booktable}

\subsection*{Review Questions}\label{08-git-every-day:review-questions}

\begin{enumerate}
\def\labelenumi{\arabic{enumi}.}
\tightlist
\item
  Name the three places your work lives, and the commands that move it.
\item
  What's the difference between \texttt{git\ diff} and \texttt{git\ diff\ -\/-staged}?
\item
  Why ignore \texttt{.env}, and why isn't that enough once it's committed?
\item
  When does a merge fast-forward, and when does it create a merge commit?
\item
  You pushed a bad commit an hour ago. Should you use \texttt{reset} or \texttt{revert}? Why?
\end{enumerate}

\begin{understandbox}

\textbf{You understand this when} you can make a change, review it with \texttt{git\ diff}, commit it with a clear message, and undo it the safe way.

\end{understandbox}
